\subsection{用化学对象表达一个两步反应机理}
\label{sec:ni1-chemical-objects}

Guided 的 NI1 运行研究甲胺对甲醛的加成，以及随后脱水形成亚胺—水复合物的过程（图~\ref{fig:ni1-chemical-objects}）。物种对象保存优化后的反应物复合物、半缩胺中间体和产物复合物，两个步骤对象分别通过一阶鞍点连接相邻物种。图中的结构和电子能全部来自该次运行自身的 B3LYP-D3(BJ)/def2-SVP 气相计算，电荷为0，自旋多重度为1。

加成鞍点具有唯一的虚频 $-1481.06$ cm$^{-1}$，相对反应物复合物的电子能垒为28.12 kcal/mol。两个方向的短 IRC 外帧分别经几何优化和频率计算到达反应物复合物与半缩胺中间体，两者均无虚频。脱水鞍点具有唯一的虚频 $-1796.84$ cm$^{-1}$，两个方向同样得到半缩胺和亚胺—水复合物。每次 IRC 均在两个方向各追踪24个点后达到点数上限，随后用独立端点计算确定相应极小值。

脱水能垒还说明一个物种对象为何需要保存多个几何构象。该步骤实际 IRC 一侧到达的半缩胺极小值，比运行中找到的最低半缩胺构象高0.71 kcal/mol。以局部 IRC 端点为参考的能垒为57.21 kcal/mol，以最低构象为参考则为57.92 kcal/mol。步骤对象因此同时保存真实路径端点和计算能垒时采用的稳定参考构象。两个构象间的互变路径仍待计算。SI 第S4节给出绝对能量、原生计算位置和可复现的绘图输入。

